\section{Equações de Lagrange e estrutura das forças de inércia}
\label{sec_dinAcoplada_equacoes_Lagrange}


As equações de Lagrange para o sistema espaçonave--\gls{mr}, em termos das coordenadas generalizadas $\vec y\,$ e de suas derivadas, podem ser escritas componente a componente conforme \cite{Wilde2018}:

\begin{equation}
\frac{d}{dt}\!\left(\frac{\partial T}{\partial \dot y_k}\right)
-\frac{\partial T}{\partial y_k}
=
\tau_k,
\qquad
k = 1,\dots,n_g.
\label{eq_lagrange_componentes}
\end{equation}

O termo $T$ é a energia cinética total, $\tau_k$ é a força generalizada associada à coordenada $y_k$ e $n_g$ é o número de graus de liberdade considerados na dinâmica acoplada (espaçonave e \gls{mr}).
Partindo da forma quadrática da energia cinética em \eqref{eq_T_quadratica_Mq1},
em que $M(\vec y\,)$ é a matriz de inércia generalizada do conjunto espaçonave--\gls{mr},
reescreve-se $T$ na forma escalar:

\begin{equation}
T(\vec y,\dot{\vec y\, })
=
\frac{1}{2}
\sum_{k=1}^{n_g}\sum_{j=1}^{n_g}
m_{kj}(\vec y\, )\,\dot y_k \dot y_j.
\label{eq_T_expandida_mkj}
\end{equation}


Sendo assim, a derivada parcial de $T$ em relação à velocidade generalizada
$\dot y_k$ é obtida de forma explícita, e ao reescrever $T$ na forma:

\begin{equation}
T(\vec y,\dot{\vec y\, })
=
\frac{1}{2}
\sum_{j=1}^{n_g}\sum_{l=1}^{n_g}
m_{jl}(\vec y)\,\dot y_j \dot y_l.
\end{equation}

Tem-se:

\begin{equation}
\begin{aligned}
\frac{\partial T}{\partial \dot y_k}
&=
\frac{\partial}{\partial \dot y_k}
\biggl[
\frac{1}{2}
\sum_{j=1}^{n_g}\sum_{l=1}^{n_g}
m_{jl}(\vec y)\,\dot y_j \dot y_l
\biggr]
\\
&=
\frac{1}{2}
\sum_{l=1}^{n_g}
m_{kl}(\vec y)\,\dot y_l
+
\frac{1}{2}
\sum_{j=1}^{n_g}
m_{jk}(\vec y)\,\dot y_j
\\
&=
\sum_{j=1}^{n_g}
m_{kj}(\vec y)\,\dot y_j
\end{aligned}.
\label{eq_dT_dqdot_intermediaria}
\end{equation}

O termo, no último passo, utiliza-se a simetria da matriz de inércia generalizada, $m_{jk}(\vec y) = m_{kj}(\vec y)$ \cite{Geradin2004}.
Em seguida, ao aplicar a derivada temporal total a \eqref{eq_dT_dqdot_intermediaria}, obtém-se:

\begin{equation}
\begin{aligned}
\frac{d}{dt}\!\left(\frac{\partial T}{\partial \dot y_k}\right)
&=
\frac{d}{dt}
\biggl(
\sum_{j=1}^{n_g}
m_{kj}(\vec y)\,\dot y_j
\biggr)
\\
&=
\sum_{j=1}^{n_g}
\bigl(
\dot m_{kj}(\vec y)\,\dot y_j
+
m_{kj}(\vec y)\,\ddot y_j
\bigr).
\end{aligned}
\end{equation}

Sendo:

\begin{equation}
\dot m_{kj}(\vec y\,)
=
\sum_{l=1}^{n_g}
\frac{\partial m_{kj}(\vec y\,)}{\partial y_l}\,\dot y_l.
\end{equation}


Substituindo essa expressão na relação anterior e reorganizando os termos, obtém-se uma forma compacta para o primeiro termo à esquerda de \eqref{eq_lagrange_componentes},
associado à coordenada $y_k$ do sistema espaçonave--\gls{mr}.
O segundo termo de \eqref{eq_lagrange_componentes} envolve a derivada parcial de $T$ em relação a $y_k$.
Usando novamente a forma escalar em \eqref{eq_T_expandida_mkj}, tem-se

\begin{equation}
\begin{aligned}
\frac{\partial T}{\partial y_k}
&=
\frac{\partial}{\partial y_k}
\biggl[
\frac{1}{2}
\sum_{j=1}^{n_g}\sum_{l=1}^{n_g}
m_{jl}(\vec y)\,\dot y_j \dot y_l
\biggr]
\\
&=
\frac{1}{2}
\sum_{j=1}^{n_g}\sum_{l=1}^{n_g}
\frac{\partial m_{jl}}{\partial y_k}\,\dot y_j \dot y_l.
\end{aligned}
\end{equation}

O termo as velocidades generalizadas $\dot y_j$ e $\dot y_l$ são tratadas como constantes na derivação em relação a $y_k$.
Essa expressão fornece o segundo termo cinético das equações de Lagrange em \eqref{eq_lagrange_componentes}.


%%%%%%%%%%%%%%%%%%%%%%%%%%%%%%%%%%%%%%%%%%%%%%%%%%%%%%%%%%%%%%%%%%%%%%%%%%%%%%%%%%%%%%%%%%%%%%%

Combinando as expressões obtidas para $\dfrac{d}{dt}\!\left(\dfrac{\partial T}{\partial \dot y_k}\right)$
e para $\dfrac{\partial T}{\partial y_k}$, as equações de Lagrange podem ser escritas na forma escalar:

\begin{equation}
\sum_{j=1}^{n_g}
m_{kj}(\vec y\,)\,\ddot y_j
+
\sum_{j=1}^{n_g}\sum_{l=1}^{n_g}
\bigg(
\frac{\partial m_{kj}}{\partial y_l}
-
\frac{1}{2}\frac{\partial m_{lj}}{\partial y_k}
\bigg)\,\dot y_l \dot y_j
=
\tau_k - g_k(\vec y\,),
\qquad
k = 1,\dots,n_g.
\label{eq_lagrange_componentes_expandida}
\end{equation}

O termo $y_k$ e $y_l$ são as componentes do vetor de coordenadas generalizadas
$\vec y$, e $g_k(\vec y\,)$ é a $k$-ésima componente do vetor
$\vec G(\vec y\,)$ associado às forças generalizadas derivadas do
potencial configuracional $V(\vec y\,)$, com
$g_k(\vec y\,) = \dfrac{\partial V}{\partial y_k}$ \cite{Geradin2004}.
Ao agrupar os termos proporcionais às acelerações $\ddot y_j$ na matriz
$M(\vec y\,)$, os termos quadráticos em $\dot y_j$ e $\dot y_l$ na
matriz $C(\vec y\,,\dot{\vec y\,})$ e as contribuições do potencial no
vetor $\vec G(\vec y\,)$, a expressão
\eqref{eq_lagrange_componentes_expandida} é organizada em forma
matricial \cite{Geradin2004}:

\begin{equation}
M(\vec y\,)\,\ddot{\vec y\,}
+
C(\vec y\,,\dot{\vec y\,})\,\dot{\vec y\,}
=
\vec\tau,
\label{eq_dinamica_geral_com_G}
\end{equation}

\begin{equation}
c_{kj}(\vec y\,,\dot{\vec y\,})
=
\sum_{l=1}^{n_g}
\left(
\frac{\partial m_{kj}(\vec y\,)}{\partial y_l}
-
\frac{1}{2}\frac{\partial m_{lj}(\vec y\,)}{\partial y_k}
\right)\dot y_l,
\qquad
k,j = 1,\dots,n_g.
\label{eq_coeficientes_C_geradin}
\end{equation}

O termo o termo $C(\vec y\,,\dot{\vec y\,})\,\dot{\vec y\,}$ agrupa as contribuições centrífugas e de Coriolis.
Conforme discutido na Seção~\ref{sec_dinAcoplada_energia_cinetica},
no modelo adotado para o conjunto espaçonave--\gls{mr} em ambiente de
microgravidade assume-se $V(\vec y\,) \equiv 0$, de modo que
$\vec G(\vec y\,) = \vec 0$ e a Lagrangiana reduz-se a
$L(\vec y\,,\dot{\vec y\,}) = T(\vec y\,,\dot{\vec y\,})$
\cite{wie2008space,murray1994mathematical}.
Nesse caso, as equações de movimento tomam a forma padrão:

\begin{equation}
M(\vec y\,)\,\ddot{\vec y\,}
+
C(\vec y\,,\dot{\vec y\,})\,\dot{\vec y\,}
=
\vec\tau.
\label{eq_dinamica_geral}
\end{equation}

O termo $M(\vec y\,)$ é a matriz de inércia generalizada construída a partir dos jacobianos espaciais
e dos operadores de inércia da espaçonave e dos elos,
conforme \eqref{eq_Ms_total_formula}, e $C(\vec y\,,\dot{\vec y\,})\,\dot{\vec y\,}$
representa o vetor de forças centrífugas--Coriolis \cite{featherstone2008,lynch2017modern}.
A organização das coordenadas generalizadas segue a definição
\eqref{eq_q_generalizado_espaconave_mr}, isto é:

\begin{equation}
\vec y
=
\begin{bmatrix}
\vec\eta_B \\
\vec\xi
\end{bmatrix},
\qquad
\dot{\vec y}
=
\begin{bmatrix}
\dot{\vec\eta}_B \\
\dot{\vec\xi}
\end{bmatrix}.
\label{eq_y_def_reprise}
\end{equation}


O termo $\vec\eta_B$ reúne os parâmetros de atitude da espaçonave e $\vec\xi$
contém as coordenadas articulares do manipulador.
De forma compatível, o vetor de forças generalizadas é dividido como:

\begin{equation}
\vec\tau
=
\begin{bmatrix}
\vec\tau_{\text{att}} \\
\vec\tau_{\text{juntas}}
\end{bmatrix}.
\label{eq_tau_estruturado_final}
\end{equation}

O termo $\vec\tau_{\text{att}}$ é o vetor de torques aplicados à atitude da
espaçonave, gerados pelo controlador de atitude, e
$\vec\tau_{\text{juntas}}$ é o vetor de torques aplicados nos atuadores
das juntas do \gls{mr}, gerados pelos controladores articulares.
Essa estrutura deixa explícito que, no modelo considerado, $\vec\tau$
agrupa apenas as ações de controle internas ao conjunto
espaçonave--\gls{mr}, enquanto as contribuições derivadas de potenciais
externos (como torques gravitacionais) são negligenciadas ao se adotar
$V(\vec y\,) = 0$.

%%%%%%%%%%%%%%%%%%%%%3333333333333%%%%%%%%%%%%%%%%%%%%%%%%%%%%%%%%%%%%%%%%

Sendo assim, é conveniente explicitar a Equação~\eqref{eq_dinamica_geral}
para o conjunto espaçonave--\gls{mr} por meio do particionamento,
conforme \cite{featherstone2008,Wilde2018}, da matriz de inércia e da
matriz de Coriolis:

\begin{equation}
M(\vec y\,) =
\begin{bmatrix}
M_{BB}(\vec y\,) & M_{BM}(\vec y\,) \\
M_{MB}(\vec y\,) & M_{MM}(\vec y\,)
\end{bmatrix},
\qquad
C(\vec y\,,\dot{\vec y\,}) =
\begin{bmatrix}
C_{BB}(\vec y\,,\dot{\vec y\,}) & C_{BM}(\vec y\,,\dot{\vec y\,}) \\
C_{MB}(\vec y\,,\dot{\vec y\,}) & C_{MM}(\vec y\,,\dot{\vec y\,})
\end{bmatrix}.
\label{eq_particionamento_M_C_Lagrange}
\end{equation}

O termo o índice $B$ está associado às contribuições da espaçonave física e o
índice $M$ às contribuições do manipulador robótico.
Ao aplicar esse particionamento na Equação~\eqref{eq_dinamica_geral},
a dinâmica da atitude da espaçonave pode ser escrita na forma:

\begin{equation}
M_{BB}(\vec y\,)\,\ddot{\vec y}_B
+
M_{BM}(\vec y\,)\,\ddot{\vec y}_M
+
C_{BB}(\vec y\,,\dot{\vec y\,})\,\dot{\vec y}_B
+
C_{BM}(\vec y\,,\dot{\vec y\,})\,\dot{\vec y}_M
=
\vec\tau_B.
\label{eq_dinAcoplada_espaconave_compacta}
\end{equation}

O termo $\vec y_B$ é o vetor de coordenadas generalizadas associado à
atitude da espaçonave, $\vec y_M$ reúne as coordenadas das juntas do
manipulador e $\vec\tau_B$ é o vetor de forças generalizadas aplicadas
à espaçonave.
Para explicitar o papel das reações inerciais geradas pelo
manipulador, decompõe-se o esforço generalizado da espaçonave em:

\begin{equation}
\vec\tau_B
=
\vec\tau_{\text{att}}
+
\vec\tau_{\text{reacao}}.
\label{eq_decomposicao_tauB}
\end{equation}


O termo $\vec\tau_{\text{att}}$ representa o torque de controle aplicado
pelos atuadores de atitude da espaçonave e $\vec\tau_{\text{reacao}}$
representa o torque interno transmitido ao corpo da espaçonave pela
movimentação do manipulador, isto é, a reação de acoplamento dinâmica
entre espaçonave e \gls{mr}. 